\section{补充推导：从 profiling 到 Triton kernel 的完整工作流}

\subsection{GPU hardware ledger}

\begin{longtable}{p{0.22\textwidth}p{0.22\textwidth}p{0.46\textwidth}}
\toprule
\textbf{硬件层级} & \textbf{典型容量/特征} & \textbf{优化含义} \\
\midrule
Registers & 每个 thread 私有，极快 & register 用太多会降低 occupancy。 \\
Shared memory / L1 & 每个 SM 内，快但小 & 适合 tile、row reduction、softmax 等 block 内复用。 \\
L2 cache & 全 GPU 共享缓存 & 对跨 block 的重复访问有帮助，但不可完全依赖。 \\
HBM & GPU global memory，大容量高带宽 & 仍远慢于片上存储；kernel 优化要减少 HBM 往返。 \\
Tensor cores & 专用矩阵硬件 & 需要 shape/dtype/layout 匹配，才能达到高吞吐。 \\
\bottomrule
\end{longtable}

\begin{warningbox}{occupancy 不是越高越好}
高 occupancy 能隐藏 HBM latency，但如果每个 thread 做更多 work、复用更多数据，较低 occupancy 也可能更快。正确问题不是“occupancy 是否最大”，而是瓶颈是不是因为可运行 warps 不够。
\end{warningbox}

\subsection{Benchmarking 的最小正确性条件}

\begin{lstlisting}[caption={Benchmarking checklist}]
# 1. Warm up first to avoid compilation / cache effects.
for _ in range(num_warmups):
    run()
torch.cuda.synchronize()

# 2. Use CUDA events for GPU timing.
start = torch.cuda.Event(enable_timing=True)
end = torch.cuda.Event(enable_timing=True)
start.record()
run()
end.record()
torch.cuda.synchronize()
time_ms = start.elapsed_time(end)
\end{lstlisting}

\begin{knowledgebox}{为什么第一次运行不能计入 benchmark}
第一次运行可能包含 CUDA context 初始化、kernel compilation、allocator warmup、cache miss 等开销。我们关心 steady-state kernel performance，所以需要 warmup，并且要用多次 trial 看方差。
\end{knowledgebox}

\subsection{Profiling 如何读 kernel 名}

Profiler 不只是给总时间，还暴露实际调用的 CUDA kernels。例如一个名字里可能出现：

\begin{itemize}
    \item \texttt{cutlass}：NVIDIA 线性代数 kernel library。
    \item \texttt{sm100}：对应 Blackwell/B200 架构。
    \item \texttt{f32}：输入/输出/accumulator dtype 信息。
    \item \texttt{64x64x16}：tile shape，暗示 block 内计算组织。
\end{itemize}

\begin{importantbox}{profile 的核心价值}
Benchmark 告诉你慢，profile 告诉你为什么慢。比如同样是 matmul，小 shape 可能调用 SIMT kernel，大 shape 才调用 tensor-core-friendly GEMM；同样是 GeLU，naive 写法可能拆成多个 pointwise kernels，compiled 版本可能融合成 Triton kernel。
\end{importantbox}

\subsection{GeLU fusion 的 HBM 账本}

Naive GeLU 公式是：

\[
\mathrm{GeLU}(x)\approx 0.5x\left(1+\tanh\left(0.79788456(x+0.044715x^3)\right)
\right).
\]

若每个子表达式都变成单独 kernel，可能产生多次 HBM read/write。Fused kernel 则把中间值保留在 registers 中：读一次 \(x\)，计算多个子表达式，写一次 \(y\)。

\begin{longtable}{p{0.24\textwidth}p{0.34\textwidth}p{0.32\textwidth}}
\toprule
\textbf{实现} & \textbf{kernel 形态} & \textbf{资源效果} \\
\midrule
Naive PyTorch & 多个 pointwise kernels & 多次 HBM 往返，launch overhead 多。 \\
Builtin GeLU & 库内 fused kernel & 减少中间读写。 \\
\texttt{torch.compile} & 编译生成 fused Triton kernel & 保留 Python 语义，获得接近手写 kernel 的数据移动模式。 \\
\bottomrule
\end{longtable}

\subsection{Triton GeLU kernel 逐行解释}

\begin{lstlisting}[caption={Triton GeLU kernel 核心结构}]
pid = tl.program_id(axis=0)
start = pid * BLOCK_SIZE
offsets = start + tl.arange(0, BLOCK_SIZE)
mask = offsets < num_elements
x = tl.load(x_ptr + offsets, mask=mask)
# compute gelu in registers
...
tl.store(y_ptr + offsets, y, mask=mask)
\end{lstlisting}

\begin{knowledgebox}{Triton program 的心智模型}
每个 Triton program 类似一个 thread block，负责一段 offsets。Mask 处理尾部不满 block 的元素。`tl.load` 从 global memory 读入向量，后续计算在 registers 中完成，`tl.store` 写回。这个模型比 CUDA thread-level 更粗，但足以表达很多 ML kernels。
\end{knowledgebox}

\subsection{Softmax 为什么适合 row-wise block}

Softmax 对每一行独立：先求最大值 \(m\)，再计算指数和归一化。

\[
\mathrm{softmax}(x_i)=\frac{e^{x_i-m}}{\sum_j e^{x_j-m}},\qquad m=\max_j x_j.
\]

若一行能放进一个 block，Triton 可以一次读入整行，在 block 内完成 max、exp、sum、divide，再写回输出。这样避免 naive 版本多次对同一行读写 HBM。

\begin{warningbox}{softmax 的数值稳定性}
先减去行最大值 \(m\) 不是性能技巧，而是数值稳定技巧。否则 \(e^{x_i}\) 可能 overflow。性能优化不能破坏这种稳定性约束。
\end{warningbox}

\subsection{Row sum：当一行放不进一个 block}

Row sum 是 softmax 的简化版，用来讲解“行太长”的情形。如果 \(N=4096\) 而 block size 是 1024，一行需要 4 个 tiles。每个 program 在 tile loop 中累加 partial sum，最后做 block 内 reduction。

\begin{importantbox}{Row sum 是 baby tiling}
它没有 matmul 那么复杂，但已经包含 tiling 的核心：把大数据拆成可放进 block 的块，循环加载，局部累加，最后归约。这是理解 FlashAttention 和 tiled matmul 的前置台阶。
\end{importantbox}

\subsection{Tiled matmul + ReLU 的完整账本}

Naive matmul 对每个输出元素 \(C_{mn}\) 遍历 \(k\)，不断从 HBM 读 \(A_{mk}\) 和 \(B_{kn}\)。相邻输出元素会重复用到同一行 A 或同一列 B。如果直接从 HBM 重复读，arithmetic intensity 低。

Tiling 做法是：

\begin{enumerate}
    \item 选择一个 \(C\) 的 output tile。
    \item 依次加载对应的 A tile 和 B tile 到 shared memory。
    \item 用 \texttt{tl.dot} 在 tile 上累积 partial sums。
    \item 在写回 C 之前顺手做 ReLU，完成 fusion。
\end{enumerate}

\begin{knowledgebox}{为什么 tile size 决定性能}
Tile 太小，数据复用不够；tile 太大，shared memory/register 不够、occupancy 下降。高性能 GEMM kernel 的核心就是在 tile shape、register pressure、memory coalescing、tensor core layout 之间找平衡。
\end{knowledgebox}

\subsection{本章小结}

这一节把源码中的 kernel examples 连接成完整工作流：GeLU 展示 pointwise fusion，softmax 展示 row-wise reduction，row sum 展示跨 tile accumulation，matmul+ReLU 展示 shared-memory tiling 和 activation fusion。四者共同说明：Triton 的价值在于让你以 block 为单位显式组织数据移动。



